\documentclass[11pt,a4paper]{article}
\usepackage{curso}
\begin{document}
\portada{Lección}{5.7}{Abstracciones de hardware}

\noindent\textbf{Unidad 5.}\quad\textbf{Objetivos:} el problema de conectividad $N^2$ y el grafo padre fijo; embebido de menores con cadenas y elección de
su fuerza; costo de embeber un modelo denso en un grafo de grado 16.
\textbf{Prerrequisitos:} 4.1, 5.3.

\section{El problema \texorpdfstring{$N^2$}{N2}}
Un Ising completo tiene $|E|=N(N-1)/2$ acoplamientos; un grafo de grado fijo $d$, $|E_d|=Nd/2$. Para $N=10^4$ y $d=16$: $8\times10^4$ frente a $5\times10^7$.
Además, en un chip los cables deben ser cortos: los vecinos de cada pbit están cerca en el plano. Ese grafo fijo y disperso es el \emph{grafo padre}.

\section{Embebido de menores}
\begin{definicion}[ (cadena)]
Una variable lógica de grado $k>d$ se representa con $c$ pbits físicos unidos por acoplamientos ferromagnéticos fuertes $J_c$. Cada copia aporta vecinos
libres; en una cadena, las copias internas usan 2 aristas para sus vecinas de cadena:
\[c\ \ge\ \Big\lceil\frac{k-2}{d-2}\Big\rceil\qquad(k>d).\]
\end{definicion}
\begin{ejemplo}
$d=16$: $k=15\Rightarrow c=1$; $k=31\Rightarrow c=3$; $k=99\Rightarrow c=7$; $k=999\Rightarrow c=72$.
\end{ejemplo}

\section{Fuerza de la cadena}
Para dos copias aisladas con energía $-J_cs_1s_2$, la probabilidad de que discrepen (cadena rota) es
\[P_{\text{rota}}=\frac{2e^{-\beta J_c}}{2e^{\beta J_c}+2e^{-\beta J_c}}=\frac{1}{1+e^{2\beta J_c}}.\]
\begin{center}\begin{tabular}{lccc}\toprule
$\beta J_c$ & 1 & 2 & 3\\\midrule
exacta & 0,119 & 0,018 & 0,0025\\
Gibbs, $2\times10^5$ barridos \evS & 0,119 & 0,018 & 0,0027\\\bottomrule\end{tabular}\end{center}
Una cadena demasiado fuerte se congela: voltear la variable lógica exige voltear todas las copias y la mezcla se vuelve lenta. Hay que equilibrar.
Al decodificar, las cadenas rotas se resuelven por voto mayoritario o se descartan.

\section{Costo en el Z1}
Z1: $269\,568$ pbits $\times16/2\approx2{,}2\times10^6$ acoplamientos. Un modelo disperso de grado 16 usa $269\,568$ variables; uno denso de mil variables
necesita 72 copias por variable, es decir, el chip alcanzaría para $\approx3\,744$ variables sin contar el ruteo. Conclusión: diseñar modelos dispersos y
locales desde el principio (Unidades 9--10). El coloreo del grafo padre fija las fases de reloj por barrido (lección 4.3).

\begin{aplicacion}
Laboratorio L8: construir un grafo de grado 16, colorearlo y embeber un modelo denso pequeño; medir rupturas en función de $\beta J_c$.
\end{aplicacion}

\begin{error}
\begin{itemize}[nosep]
\item Suponer que cualquier modelo corre directamente. \quad Ignorar las cadenas rotas al decodificar.
\item Cadenas tan fuertes que la mezcla se congela. \quad Contar copias sin contar el ruteo.
\end{itemize}
\end{error}

\begin{thebibliography}{9}\small
\bibitem[22]{z1} Extropic, From One to One Billion: Torx, Thermalizers, and Z1, 2026.
\bibitem[19]{aadit} N. A. Aadit et al., Massively parallel probabilistic computing with sparse Ising machines, 2022.
\bibitem{choi} V. Choi, Minor-embedding in adiabatic quantum computation, Quantum Inf. Process. 7, 193 (2008).
\end{thebibliography}
\end{document}
